% ЗАКЛЮЧЕНИЕ, без номер.
\section*{ЗАКЛЮЧЕНИЕ}
\label{sec:conclusion}
\addcontentsline{toc}{section}{ЗАКЛЮЧЕНИЕ}

В рамките на курсовия проект е разработена система за наблюдение на показатели в реално
време, изградена от сървър на C++20~\cite{iso_cpp20}, който сам обслужва
HTTP/1.1~\cite{rfc9112} и WebSocket~\cite{rfc6455}, и от клиент в браузъра без приложна
рамка. Клиентската част стъпва пряко върху обектния модел на документа~\cite{dom_ls},
модела на събитията, структурната графика~\cite{svg2} и разположението с
CSS Grid~\cite{css_grid1}, а обменът е реализиран в два формата~\cite{xml10,rfc8259}, за
да може разликата между тях да бъде измерена.

\textbf{Предимства на решението.} Отсъствието на приложна рамка прави всяка стъпка от
обработката видима в изходния текст, а сървър, който сам обработва протокола, дава пълен
контрол над момента на изпращане, което е предпоставка за измерването на закъснението.
Разделянето на сериализатора от източника на показатели позволява смяна на формата по
време на работа, без промяна в останалата част от системата.

\textbf{Недостатъци и ограничения.} Слоят за свързване на данни е написан ръчно и
изисква явно описание на всяка зависимост, което при разрастване на интерфейса става
източник на пропуски. Състоянието се пази в паметта на процеса и се губи при рестарт.
Няма удостоверяване на потребители. Тези ограничения са заявени в
Раздел~\ref{sec:design} преди измерването, а не открити след него.

\textbf{Инженерна трактовка на резултатите.} \TODO{извод от измерените стойности: какво
означават те за избора между двата канала и между двата формата в задача от този вид}

\textbf{Насоки за по-нататъшна работа.} Естествените продължения са съхранение на
исторически редове извън паметта на процеса, обединяване на няколко източника на
показатели, удостоверяване на потребители и разширяване на слоя за свързване на данни с
автоматично проследяване на зависимостите.
